\section{Three boundaries}
\label{sec:boundaries}

Figure~\ref{fig:boundaries} places the boundaries relative to the protocol. We
state each one as a claim about the MR family, name the evidence that carries it,
and say what would falsify it.

\begin{figure}[t]
\centering
\includegraphics[width=0.98\linewidth]{fig7_boundaries.pdf}
\caption{The three boundaries against the protocol of Figure~\ref{fig:protocol}.
Boundary~1 is the deterministic false proof produced by the adjudication as
designed; boundary~2 is the emptiness of the obvious repair---once a readable
oracle exists, the metamorphic step is not empirically necessary; boundary~3 is the conditional
nature of the specificity figure computed over the reachability model.}
\label{fig:boundaries}
\end{figure}

\subsection{Boundary 1 (RQ1): the adjudication as designed produces deterministic false proofs}

\emph{Taken as published---raw alarms, or alarms filtered by bookkeeping
exceptions---the adjudication reports a legitimate access granted out of band
after the crawl as an already-established violation, and it does so in one
direction.}

The evidence is twofold. On the constructed system, a legitimate share created
out of band yields \ResObsOobVtwoFalseProof{} false proofs out of
\ResObsOobFalseAlarm{} alarms, a rate of \ResObsOobVtwoFpRate{}, against
\ResObsFwVtwoFalseProof{} when the same share was created through the framework's
facility. On Gitea, the bookkeeping adjudication produces
\ResRsMrBookkeepingFp{} false proofs on \ResRsAuthorizedCells{} authorized cells
(\ResRsMrBookkeepingFpRate{}), and the same method under a live reading of the
reachability model produces \ResRsMrBookkeepingAlarmsLive{} alarms---so the
error is not a property of the system, it is a property of the reading.

We stress the \emph{deterministic direction}, because the original literature
attributes its residual false alarms partly to asynchronous behaviour. Here the
direction is forced. The predicate has exactly two representations available
---``the subject cannot reach this object'' and ``the subject can reach this
object''---and a grant created outside the crawled interface has no
representation other than the first. The error is therefore systematic, which is
why the false-proof count is not merely non-zero but equal to the alarm count on
the out-of-band scenario. Section~\ref{sec:threats} records why our
approximation of the exception semantics cannot turn this into a measurement of
the original engine's behaviour.

This boundary is not our discovery. The predecessor design attributes its own
residual false alarms to the crawler failing to traverse ``all URLs for all
users''~\cite{Mai2019Metamorphic}. What we add is the identification of the
failure as a function of \emph{when} the model was read rather than \emph{which}
objects it contained, which is what makes the repair of Boundary~2 available---
and, as we now show, empty.

\subsection{Boundary 2 (RQ2): once a sound, readable oracle exists, the metamorphic step is not empirically necessary}

\emph{If the false proofs of Boundary~1 are removed by consulting a readable
authorization introspection surface, then it is the oracle that carries the
work: replacing the entire metamorphic step with a single ordinary request
changes nothing measurable.}

The evidence is the ablation of Section~\ref{sec:synthetic-ablation}.
\texttt{direct\_observe}---an ordinary request that succeeds, checked against
the same readable surface---attains false proofs of \ResAbDirectObserveFP{}
against v3's \ResAbVThreeFP{}, and recall \ResAbDirectObserveRec{} against v3's
\ResAbVThreeRec{}. The Gitea experiment reproduces the ordering on a real
system, with \texttt{direct\_observe} emitting \ResRsDirectObserveAlarms{}
alarms and \ResRsDirectObserveFp{} false proofs against v3's
\ResRsVThreeFp{}.

Three qualifications are part of the claim, not disclaimers appended to it.

The first is that the false-proof comparison is a \emph{tie at zero}, not a win.
We state it as a tie because that is what it is: once either method can read the
oracle, both stop producing false proofs, and the surviving difference is on
recall alone. The pair of numbers in the previous paragraph must therefore be
read as ``no measurable change'', and a reader who takes \ResAbDirectObserveFP{}
against \ResAbVThreeFP{} as an advantage has misread us.

The second is the honest arithmetic on that surviving difference. The recall
difference is \ResAbRecallDiff{} with 95\% CI
$[\ResAbRecallDiffLo{}, \ResAbRecallDiffHi{}]$, which includes zero, so
\texttt{direct\_observe} is \emph{not} demonstrated to be better. It is
demonstrated to be \emph{not worse}, and that is already decisive, because for
the metamorphic step to carry work it would have to be better at something the
plain request cannot do. It is not: on the \ResAbDisagreementN{} records where
the two methods differ, the MR-based method is the one that fails, its
reachability precondition blocking a detection the plain request makes.

The third is that the control is \emph{constructed}, not discovered. The two
observables \texttt{direct\_observe} consumes---the request succeeded, and the
authorization surface reports no entitlement---are definitionally the same
quantities that define the ground truth. The comparison is therefore a
statement about the definitions, and it must be read as one: we do not claim a
formal non-inferiority result, because no margin was fixed in advance; when the
oracle is readable, the metamorphic relation has nothing left to compute. This is also why
we do not present it as a competing method.

\subsection{Boundary 3 (RQ3): the reported specificity rests on an undeclared premise and a low-resolution denominator}

\emph{The reported specificity of \ResLitReportedSpecPct{}\% is a statement about
a corpus of generated inputs under the assumption that the reachability model
coincides with the true authorization state. It cannot detect a failure that is
uniform across a scenario class, and that assumption is not stated.}

There are two distinct issues and we keep them apart.

\paragraph{The premise} The adjudication of an authorization MR is a function of
the reachability model (Section~\ref{sec:background}). A specificity figure
computed over that adjudication is therefore conditional on the model being
right. Where the model and the true authorization state diverge---exactly the
out-of-band case of Boundary~1---the figure is not measuring the tool's
agreement with authorization, it is measuring its agreement with its own
crawler. None of this is concealed by the authors: they say the role-to-resource
mapping is ``not always available''~\cite{Chaleshtari2022Metamorphic}, and their
predecessor names the crawler as the source of its false
alarms~\cite{Mai2019Metamorphic}. Our point is that the number is reported
without that condition attached.

\paragraph{The denominator} The published definition phrases specificity as the
fraction of \emph{generated inputs} that do not lead to a false
positive---\ResLitMaiFalseAlarms{} of \ResLitMaiGenInputs{} in the predecessor's
own accounting, a figure whose own text attributes the residual cases to crawler
coverage. A denominator in generated inputs is divided by every request the
generator emits. A failure that removes an entire scenario class from correct
handling contributes only as many false positives as that class emits inputs, so
it is diluted by the rest of the run. In our approximation
(Table~\ref{tab:spec}) the configuration in which a whole class of legitimate
access is misreported still reports a specificity of \ResSpecRowOneSpec{}.
The metric is not computed incorrectly; it is \emph{low-resolution with respect
to this failure mode}, and that is a property of the choice of denominator, not
of the arithmetic.

\paragraph{What we are not saying} We are not saying the reported figure is
false, and we are not saying the reported false positives were misdiagnosed. The
published attribution for those is asynchronous action, which is a different
cause from the one we describe. The claim is narrower: a specificity in
input-units, computed against a reachability model, cannot discharge the premise
that model embeds, and cannot see a class-uniform failure.

\subsection{Implementation-level confirmation}
\label{sec:impl}

An earlier draft of this argument rested on the published \emph{definitions} of
the two predicates. Definitions can be read generously. We therefore obtained the
implementation---the MST/SMRL library that realises \mbox{MST-wi} is
public~\cite{Chaleshtari2022Metamorphic}---and checked the predicates at source.
The library ships \ResArtMrTotal{} concrete metamorphic relations, of which
\ResArtAuthzMrTotal{} are authorization relations;
\ResArtAuthzGatedEither{} of those gate on one of the two predicates we analyse
(\ResArtAuthzGatedCrtg{} on \texttt{cannotReachThroughGUI},
\ResArtAuthzGatedUcrc{} on \texttt{userCanRetrieveContent}), and
\ResArtAuthzGatedNeither{} gate on neither.

Three facts emerged that the published text does not state.

\paragraph{The freeze is hard-coded, not incidental}
\texttt{cannotReachThroughGUI} resolves to a set-membership test: for the current
account it collects the URLs recorded in the \emph{crawl} inputs and asks whether
the URL under test is absent from that set. It issues no request.
\ifnum\ResArtFactFOneurlsetmemoizedneverinvalidated=1
The per-account set is then cached in a field that is written exactly once and
never invalidated anywhere in the library,\else
(The library differs from our reading here.)\fi{}
so the snapshot is not an artefact of when the crawl happened to run; it is
enforced by the implementation. \texttt{userCanRetrieveContent} is the same in
kind: it compares against an output store read from a filesystem directory
produced by the crawl, and that store too is loaded once and reused.

\paragraph{The clause that could have caught it is inert} The reachability test
has a second clause: the URL must also be absent from the set recorded
\emph{across all users}, which is the clause that would prevent an alarm on an
object many accounts can reach. That set is built by scanning every recorded
sequence and stopping at the first login action; every record in our snapshot
begins with a login action, so the set is empty and the clause never fires. This
is not the same as the snapshot merely \emph{having no second opinion}: the
library \emph{has} one and it cannot report, so the per-account record alone
decides the verdict.

\paragraph{Missing data defaults to an alarm}
\ifnum\ResArtFactFTwofailopenwhencrawlinputempty=1
The predicate returns \texttt{true}---``the subject cannot reach this
object''---when the crawl input list is empty, and again when either argument is
null.\else
(The library differs from our reading here.)\fi{}
An empty crawl therefore \emph{satisfies} the MR's precondition instead of
abstaining. This is the mechanism behind the deterministic direction we reported
in Boundary~1: that failure is not noise, it is a default.

\paragraph{The role relation is configuration, not state}
\ifnum\ResArtFactFFivesupervisorfromstaticconfig=1
The guard on the second MR family is a supervisor relation read from the
system's static configuration file, and it reports \texttt{true} when the two
accounts are the same.\else
(The library differs from our reading here.)\fi{}
Attributes of this kind are fixed at configuration time and cannot reflect an
authorization granted afterwards.

What this removes is a specific objection---that we inferred the predicates'
semantics from prose. The semantics are in the code. It does not by itself
establish that the code \emph{behaves} as read, so we also executed it.

\subsection{Executing the published relations}
\label{sec:native-run}

The library ships the relations but not the execution environment: no system
configuration, no crawl snapshot, no subject system. We supplied those, and ran
the relations themselves---source unmodified, with the build configuration
(\texttt{pom.xml}) and an external harness/driver added---against Gitea
\NtGiteaVersion{}.

\paragraph{What ran} \NtMrExecuted{} of the \NtMrTotal{} authorization relations
executed. The engine's own input split partitions the input view into
\NtInputView{} parts, one per crawl record, and each part is evaluated against
the \NtUserCount{} accounts it derives from those records' login actions, so
every relation is exercised over the full record--account product.
\NtMrFired{} of the \NtMrExecuted{} relations returned at least one failure,
\NtAlarms{} alarms in total. Those counts are the engine's own
\texttt{getFailures()} return value, not our adjudication of it.
Table~\ref{tab:native} lists them per relation.

\paragraph{What dominates the wall time} Almost none of the run's \NtElapsedS{}
seconds is analysis. Each of the \NtIterations{} leaf evaluations invokes the
engine's own browser-cleanup hook, which on Windows spawns two process-kill
commands and waits on each whether or not a browser is present; relations that
issue no request at all still consume minutes. This is a cost of the execution
environment the artifact does not ship, and it bounds how much of the catalog a
reproduction can exercise.

\paragraph{The load-bearing cell} Take the topology of
Section~\ref{sec:realsystem}: the subject \NtKeyGrantee{} holds a legitimate
entitlement to \NtKeyRepo{}, granted out of band, after the crawl. On that cell
the published predicate \texttt{cannotReachThroughGUI} answers ``cannot
reach''---which is the answer it is \emph{defined} to give, because the
subject's crawl record contains no entry for the repository. Across the run the
predicate was invoked \NtCrtgCalls{} times and returned that answer \NtCrtgTrue{}
times (\NtCrtgTruePct{}\%), of which \NtKeyCrtgTrue{} fall on this cell.
\NtKeyMrFired{} relations (\NtKeyMrNames{}) report a failure there. The alarm is
therefore a false proof produced by the published relations, under the published
predicate, on a real system: the subject's request succeeds and returns the same
bytes as the owner's, and the family's guarded comparison---gate on the
predicate, then compare the credentials-changed observation with the original
(Section~\ref{sec:impl})---reads an unchanged observation as a violation. The
access is authorized. Byte-equality is not inferred here from the lengths the run
happened to record: the \NtByteCells{} granted cells were re-probed afterwards, and
in each case the owner's and the grantee's response to the same URL hash
identically while a third account is refused.

The alarms are not confined to that one cell, and they go nowhere else: the
engine emitted \NtAlarms{} of them, on \NtAlarmCells{} cells, every one of them
an object to which the subject holds a grant made after the crawl\ifnum\NtAlarmsOffGrant=0,
and none on an object the subject is not entitled to.\else, and
\NtAlarmsOffGrant{} on an object the subject is not entitled to.\fi{}

\paragraph{What was and was not changed} We did not modify the library.
\texttt{git status} on the upstream tree reports exactly \NtUpstreamModified{}
changed file, \texttt{pom.xml}, and \NtUpstreamJavaModified{} changed Java
sources; the diff adds one \texttt{<source>} entry declaring the relation
directory as a build source and removes nothing. Two classes were \emph{added}: a
provider and a driver, \NtHarnessLines{} lines together. The provider reproduces
the session semantics the original replay relies on---the identity established by
a login action carries to the actions that follow it, rather than being read off
each action separately---because the credential-substitution helper rewrites only
the credential-bearing action and leaves the rest carrying the original account.
Getting this wrong does not merely add noise: it makes a credentials-changed input
fetch under the original identity, so the observation is unchanged and the
relation reports a violation on cells where the subject's own objects are at
stake. We record it because it is the kind of substitution defect that a
false-proof rate cannot distinguish from the failure being studied. The
substitution is otherwise confined to the observation channel, which the
two-stage architecture makes separable: the relations read a crawl snapshot that
the driver loads from disk, and \texttt{cannotReachThroughGUI} resolves to
set-membership over those loaded records and issues no request of its own. Every
record in that snapshot is re-read after the crawl as its own user---\NtSeedVerified{}
access points in all, each one resolving---which is necessary and not sufficient
for a crawl record, but it rules out the failure mode that mattered here: a record
no crawl could have produced. Our channel therefore cannot reach the predicate,
which is why the cell above is attributable to the definition rather than to the
harness.

Three limits belong to the claim. The content predicate
\texttt{userCanRetrieveContent} was invoked \NtUcrcCalls{} times, so this run
does not exercise it; and because that predicate is fed the observed output, our
channel \emph{would} propagate into it if it were---we report the count rather
than assert a clean separation we did not test. The supervisor guard
\texttt{isSupervisorOf} was invoked \NtSupCalls{} times
(\NtSupTrue{} true). Nor do all the relations gated on that predicate fire, and
the reasons are inside their own guards rather than in the system: one compares
the redirect URL a navigation produces, and our observations are direct responses
with no navigation, so that attribute is empty and it passes; another requires a
six-input composition our view does not furnish, so its guard is false and it
passes. Nor does every relation that runs produce a verdict: \NtMrErrored{}
aborts with no request issued at all---the engine's own reporting path
dereferences an input slot its guard never bound---and that is the one missing
from the executed count, shown as \texttt{error} in Table~\ref{tab:native}. We therefore read the alarms as
originating in the gate, not as a count of how many published relations a full
replay would flag. And Gitea is correct, so the run can only exhibit false
proofs: it cannot measure detection.

\begin{table}[t]
\centering
\caption{The published authorization relations (source unmodified; the build
configuration and a replay harness were added) executed against
Gitea \NtGiteaVersion{}. Alarms are the engine's own return value. Requests
counts observations actually issued, \NtHttpCalls{} in total: a relation
evaluated on an output already observed earlier in the same run issues none, so
requests are attributed to whichever relation observed them first.}
\label{tab:native}
\small
% 由 analysis/ingest_native_engine.py 自动生成 —— 请勿手改
% 逐条 MR 结果：告警数由引擎自己 getFailures() 返回，非本工作判读
\begin{tabular}{lrr}
\toprule
MR & alarms & requests \\
\midrule
OTG\_AUTHZ\_001 & 0 & 0 \\
OTG\_AUTHZ\_001b & 0 & 0 \\
OTG\_AUTHZ\_001b2 & 0 & 0 \\
OTG\_AUTHZ\_002 & 3 & 270 \\
OTG\_AUTHZ\_002a & 0 & 0 \\
OTG\_AUTHZ\_002b & 0 & 0 \\
OTG\_AUTHZ\_002c & 3 & 0 \\
OTG\_AUTHZ\_002d & 0 & 0 \\
OTG\_AUTHZ\_002e & \multicolumn{2}{c}{error} \\
OTG\_AUTHZ\_003 & 0 & 0 \\
OTG\_AUTHZ\_004 & 0 & 0 \\
\bottomrule
\end{tabular}

\end{table}

\subsection{The admissibility condition, and why it is the real ceiling}
\label{sec:admissibility}

The results above suggest the only repair available inside this design. An
adjudicator may promote a verdict to \texttt{violation-proven} on the strength of
a readable authorization introspection surface only when \emph{both} of the
following hold:

\begin{enumerate}
\item the application exposes a readable introspection surface, \emph{and}
\item one can give a reason that this surface is not of the same kind as the
      defect under test---that is, that the surface is not itself a
      reachability or GUI-derived artefact whose blind spot is correlated with
      the blindness being tested for.
\end{enumerate}

The second clause is the one that matters, and it is not automatable. It is a
judgement about the provenance of a piece of evidence, and it must be made per
system. The cross-vendor replication of Section~\ref{sec:gogs} makes the clause
concrete: on Gogs \GgVersion{} the only readable surface fails it, on GitLab
\GlVersion{} the readable membership surface fails it on every channel that
does not itself confer membership (it resolves inherited membership, and is
blind to public visibility and site administration), and on both systems the
oracle-gated adjudicators convert those blind spots into \GgFpVthree{} (Gogs)
and \GlFpVthree{} (GitLab) false proofs for v3, where the sound oracle of
Gitea produced none. This is why we describe the outcome as a boundary rather
than as a
mechanism: where the condition holds, the honest method is one that consumes the
oracle directly and drops the relation (Boundary~2); where it does not, the
honest method abstains, as v3 does on the unreadable stratum, where it keeps
\ResIntroOffVthreeKeepRate{} of the detections the exception-aware method keeps
and abstains \ResIntroOffAbstain{} times. We note that in that stratum the
aggregate table would have shown v3 as the best method in the study. It is not
better; it has stopped answering.

\subsection{What would falsify each boundary}

\begin{itemize}
\item \textbf{Boundary~1.} A demonstration that the predicates are defined over
      live reachability rather than recorded reachability, which would remove the
      time-dependence entirely and make the out-of-band case unremarkable.
\item \textbf{Boundary~2.} An authorization MR whose adjudication requires an
      observable that a succeeding request plus a readable oracle cannot
      produce---for instance, evidence about a subject's entitlements that the
      subject itself cannot read but the test can. We do not have such an MR; we
      have two, of the same family (Section~\ref{sec:threats}).
\item \textbf{Boundary~3.} A specificity reported in scenario units, or with the
      reachability premise made explicit, that still fails to separate the
      configurations of Table~\ref{tab:spec}.
\end{itemize}
